
\documentclass{amsart}

\usepackage{fullpage}
\usepackage[british]{babel}
\usepackage[dvipsnames]{xcolor}

% For checking overfull boxes -- Comment in final version
\overfullrule=5pt

\newcommand{\todo}[1]{{\color{red} \sf TODO: [#1]}}
\newcommand{\rs}[1]{{\color{violet} \sf RS: [#1]}} 
\newcommand{\tb}[1]{{\color{teal} \sf TB: [#1]}}
\newcommand{\pap}[1]{{\color{blue} \sf PAP: [#1]}}
\newcommand{\cm}[1]{{\color{brown} \sf CM: [#1]}}
\newcommand{\ajy}[1]{{\color{orange} \sf AJY: [#1]}}

\title{What everyone should know about scheduling computations}
\author{Papp, P. A. \and Yzelman, A. N. \and Steiner, R. S. \and ...}

\begin{document}

\begin{abstract}
Scheduling tasks on compute systems of any scale has become ubiqutous,
and is no longer a specialist topic in computer science and applied mathematics.
Over the past few years, Huawei has endeavoured to further develop the theoretical foundations of scheduling,
crucially from the perspective that in modern compute systems,
data movement is by far more costly than actual computation.
%while furthermore system architecture has become increasingly hierarchical.
This industry-wide trend has wide-ranging practical implications--
not only in areas ranging from parallel system to algorithm design,
but also on software, communication layers, and run-time systems.
The article before you summarises our recent advances,
outlines their practical implications,
and translates these to concrete advice on parallel system and software design.
\end{abstract}

\maketitle

\section{Introduction}

scheduling is everywhere: CPUs, NPUs, SOCs

scheduling becomes more intricate at scale:
\begin{itemize}
	\item multi-die,
	\item multi-device,
	\item cloud \& supercomputer scale
\end{itemize}

computation as a (Hyper)DAG

scheduling is not partitioning

\subsection{Computation model}
\label{sec:cost-models}

PRAM is not sufficient, must cost communication

note that ideally, furthermore: hierarchical
\begin{itemize}
	\item hierarchical processor designs, hierarchical nodes, hierarchical networks
\end{itemize}
latencies matter (can cite Anastasiadis if ready)

\subsection{Applications}

a summary of where scheduling is applied at present
\begin{itemize}
	\item run-time systems, such as OpenMP;
	\item communication layers, such as MPI, GASNET, others (PRAM/PGAS?);
	\item within compilers (to various degrees)
	\item offline system design (e.g., Wireless)
	\item not meant as an exhaustive list of applications, but rather to index the quite different domains of applicability
\end{itemize}


\section{Scheduling is hard}

It has been long known that scheduling is a particularly challenging problem. The complexity of problems is computer science is often defined via the concept of NP-hardness. Intuitively speaking, an NP-hard problem is one that is very unlikely to have an algorithm that can always find the optimal solution efficiently (i.e.\ with a running time that is polynomial in the size of the problem). In particular, if there was an efficient algorithm for any of the NP-hard problems, then all the other NP-hard problems could also be solved efficiently by reducing these problems to each other. There are thousand of problems that are known to be NP-hard, and many of these have been extensively studied by the community for decades; as such, when a problem is NP-hard, it is generally considered very unlikely that a polynomial-time algorithm exists for the problem.

Besides NP-hardness, there are other concepts that show that the problem is even more challenging; for instance, inapproximability means that it is even NP-hard to find a solution that is close to the optimal value, whereas parameterized complexity studies whether the problem remains hard if some parameters of the input are fixed to a small value.

With the above tools, it has been shown that scheduling is a hard problem even in very simple cases: with unit execution time for all tasks and unit communication delays between the processors, finding the optimal schedule is still NP-hard~\cite{rayward1987uet}. In more realistic settings, like the BSP model, the problem is even more difficult. For instance, the optimal solution is still NP-hard when considering very special classes of DAGs, such as DAGs with only two layers, or in-trees where each vertex has at most one outgoing edge. Moreover, in general DAGs, the optimal cost is even NP-hard to efficiently approximate beyond a specific ratio $(1+\delta)$, for some $\delta>0$; in official terms, the problem does not allow for a so-called Polynomial-Time Approximation Scheme. In fact, in BSP scheduling, even if the assignment of nodes to processors and supersteps is already fixed, the subproblem of scheduling just the necessary communication steps into the communication phases is already NP-hard~\cite{BSP_DAG_opdas}.

On the positive side, some theoretically efficient (polynomial-time) algorithms do exists for edge cases, e.g.\ when the DAG only consists of independent chains and the number of processors is small, but these algorithms are also very inefficient in practice~\cite{BSP_DAG_opdas}.

\todo{layman terms needed here?}

\subsection{Two-level scheduling}

NP-, APX-hardness, counter-examples, arbitrary factors away from optimum

\subsection{Vertical data movement: communication \emph{and} I/O}

Another vital aspect of scheduling computations in practice is vertical data movement. This concept of not included in the base BSP model for simplicity: BSP only considers the cost of moving data horizontally from the memory associated with one processor to memory associated with another. However, when executing real computations (even on just a single processor), one of the most important factors is that fast memory in the system (e.g.\ cache) often has very limited capacity, whereas we also have some other memory (e.g.\ RAM) with drastically larger capacity, but using this is much slower. In particular, the cost of moving data from slow memory to fast memory (or vice versa) is often the bottleneck in the execution of complex computations.

This vertical communication cost (also known as I/O cost) is most often analyzed in the so-called red-blue pebble game of Hong and Kung~\cite{HongKung}. This is essentially a single-processor scheduling problem in a two-level memory hierarchy (with e.g.\ cache and RAM), where the goal is to minimize the I/O operations while executing the computations. The red-blue pebble game has been widely used to analyze the I/O cost of algorithms and also to develop I/O lower bounds for many important computations, such as matrix multiplication, fast Fourier transform or attention in neural networks~\cite{HongKung,Saha}. However, in general DAGs, finding the optimal I/O cost is again a problem that is NP-hard to even approximate to any reasonable factor~\cite{MPP_opdas}.

For a more realistic model of paralellizing computations, one needs to combine the BSP model, which allows multiple processors and horizontal communication, with this red-blue pebble game that captures memory limitations and I/O costs for the single-processor case. Simply generalizing the red-blue pebble game to a parallel setting results in a multi-processor red-blue pebbling model~\cite{MPP_opdas}, whereas if we also incorporate vertex weights and synchronization costs from BSP, we obtain a more complex Multi-BSP (or MBSP) model~\cite{multiBSP_opdas}. These combinations offer much more realistic models for the scheduling problem, but on the other hand, finding the optimal solution becomes even more challenging. In particular, the hardness results for both BSP scheduling and red-blue pebbling carry over to these combined models~\cite{MPP_opdas}.

However, even if the optimal schedules are essentially impossible to find, the theoretical and experimental study of these models still offers valuable insights. For instance, one can analyze the difference between holistic scheduling strategy that aims to find a solution for the entire problem, and a two-stage approach that first comes up with a good parallelization strategy, and then combines this with a cache management strategy to optimize I/O costs. On the theoretical side, one can show a worst-case example where the two-stage schedule is a very large factor away from the actual optimum, even if both stages are optimal separately. On the experimental side, one can confirm that the holistic view can indeed outperform the two-stage approach by a small margin~\cite{multiBSP_opdas}.

\todo{Savage needed?}


\subsection{Partial computations}

The multi-processor red-blue pebbling or MBSP models above already capture most of the key factors of scheduling a computation. However, there are of course many further aspects that could be incorporated into the model to make it even more realistic. Here we briefly discuss one of these extensions: partial computations.

For simplicity, the original red-blue pebble game assumes that every operation is computed in one go, with all of its inputs in fast memory simultaneously. However, if the operation is e.g.\ the addition or multiplication of numerous inputs, then in practice, it can also be executed in multiple steps: we first aggregate only some of the input values, maybe even save this partial results into slow memory and load it back later, then continue with another subset of the inputs, and so on. Such a partial computing strategy might allow for a lower I/O cost in some cases, or it might even be necessary to allow a feasible solution when an operation has a very high number of inputs.

The red-blue pebble game can naturally be extended with such partial computations to allow for a more realistic model of I/O costs~\cite{partial_opdas,sobczyk2026complexity}. From a theoretical perspective, this extended model is rather similar to the original game in terms of e.g.\ inapproximability. Moreover, one of the most important tools related to red-blue pebbling is so-called S-partitions, which were widely used to derive lower bounds on matrix multiplication, fast Fourier transform, neural network attention, and many other computations~\cite{HongKung,Saha}. These S-partitions and the corresponding proofs can also be adjusted to the partial computing model, which allows to derive the same I/O lower bounds with partial computations (up to constant factors) for the example computations mentioned above.
\todo{state the concrete bounds? it would require introducing a LOT of notation...}  

However, we also emphasize that in general, partial computations can enable significantly different pebbling strategies. For instance, they can notably reduce the optimal I/O cost in sparse matrix-vector multiplications, or even by a linear factor in some synthetic examples~\cite{partial_opdas}.

\todo{Determining optimal schedules using ILP formulations at least as hard as the non-partial case (?) - needed?}


\section{Avoiding hardness}

While general scheduling is hard in a variety of realistic models,
scheduling in practice works around these limitations in a variety of ways.
We index common approaches along the following practical use cases:
\begin{enumerate}
	\item minimise horizontal data movement (communication);
	\item minimise vertical data movement (I/O); and
	\item minimise data movement through NUMA hierarchies.
\end{enumerate}
The approaches described below include proven techniques by which the high-performance computing (HPC) community has operated in the last few decades,
such as proving optimality for specific problems or specific classes of problems,
as opposed to attempting to find solutions that work in the most generic cases.
%Such problem-specific approaches have reliably led to algorithms and software that indeed achieve optimal performance in practice.
We couple outlining these state-of-the-art methodologies with novel insights from our research,
and close with practical implications to parallel frameworks for AI and beyond,
run-time system design for modern large-scale compute infrastructures,
and insights tailored to the complex case of irregular computations.

\subsection{Horizontal data movement}

Minimisation of horizontal data movement, i.e., communication minimisation, comes naturally when adopting a model of computation that assigns explicit costs to data movement and synchronisation.
One such model, the Bulk Synchronous Parallel (BSP) model, %developed during the late 80s,
comes with two modes of operation, %:
%one automatic, where programs are expressed sequentially and through hashing of program data, executable in parallel in a fully scalable fashion; and
%one manual, 
one of which is the manual mode
where programs explicitly coordinate computation and communication%, and potentially differ at different processes
~\cite{valiant90}.
In both cases, the cost model revolves not only around work-- but also supersteps and $h$-relations that together quantify the cost of data movement-- see Section~\ref{sec:cost-models} for details.
Costs are typically parametrised in the number of processors $p$ available and some problem-specific parameter $\vec{n}$, such as a problem size.
%To recall, the related system parameters are $g$, the cost of sending one word, and $l$, the cost of moving from one superstep to the next-- principially, making sure that any necessary communications are completed before continuing, and $p$, the number of processing units in the parallel system.
%Compute costs are hence expressed in terms of $g$ and $l$, while parameterised in $p$ and some problem size $n$.
%AJ: check later if we do indeed need to recall more of these definitions
The key phrase here is \emph{problem-specific}: in HPC, by far most known algorithms are manual-mode BSP, and literature is rich with examples of optimal algorithms or algorithms that operate within an optimal trade-off regime.
%By contrast, most highly scalable big-data frameworks instead operate by automatic-mode BSP, and achieve an \emph{expected} cost of $\mathcal{O}(n/p)$, where $n$ here could be taken as proportional to the number of data accesses the sequential program requires.
We give two examples, fast Fourier transformation (FFT) and matrix multiplication, and demonstrate how optimal algorithms and optimal trade-offs may be derived.

\subsubsection*{FFT}

introduce notion of immortal algorithms

\subsubsection*{Matrix--matrix multiplication}

introduce notion of polyalgorithms

\subsubsection*{Generalisability}

positive: these are tried an tested approaches. negative: needs human in the loop

novel approaches that automate usage of these insights, e.g., ALP

note that irregular problems cannot escape hardness, e.g., SpMV, and point forward to the specific subsection

\subsection{Vertical data movement}

This is a much harder problem, solvable to optimality for severely small problems in the general case
\begin{itemize}
	\item (point forward to Section on OneStopParallel)
\end{itemize}
Similar implications therefore apply here: limit to specific problems to find lower bounds, then find algorithms that attain them (e.g., GEMV)

\subsection{NUMA}

be careful with multi-level (distributed) scheduling

gets worse with more NUMA levels (recap SPAA results)

arbitrary factors away from optimal applies when a scheduler has no knowledge about what other schedulers are doing
\begin{itemize}
	\item therefore, do distributed scheduling where schedulers have a (compressed/sketched) view of other schedulers: \emph{federated scheduling}
\end{itemize}

Tools for general problem exist: recall Jain \& Zaharia, point out weaknesses and differences, point again to OSP for heuristics that work on multi-processor pebbling

\subsection{Parallel frameworks and run-time systems}

A section on general advice for parallel frameworks and run-time systems

\subsection{Irregular computations}

Index good approaches for problems that are inherently irregular (and thus cannot, to significant degree, avoid hardness)


\section{Similar problems, similar results, and solutions}

Partitioning,
streaming,
imply the same work-arounds.


\section{OneStopParallel}

Describe what OSP contains and aims to provide/support

\subsection{Optimal Scheduling}

Describe ILP-based formulations (don't give the formulations though, keep it high level)

\subsection{Heuristics}

Describe heuristics that are stand-alone and/or part of multi-level / divide-and-conquer approaches in the next section

\subsection{Multi-level}

Describe approaches suitable for cutting up large problems into smaller parts, ILP-based and otherwise

\subsection{Other features}


\section{Results}

 Summarise key results from the various publications and applications
 \begin{itemize}
	 \item EDA / SpTrsv
	 \item MoE
	 \item algorithm-specific parallelisation should be pointed out? (FFT, GEMM, Attention?, etc.)
 \end{itemize}


\section{Conclusion}

For important kernels or classes of computation, nothing beats manual work (yet:)-- lower bounds, trade-off identifications, (poly-)algorithms that achieve them, combined with engineering efforts. Getting bounds though should not be skipped (engineers need to know what to aim for)

Include also mention of OSP-as-a-Service, can deploy in application, or use for testing/exploration

\subsection{Future Work}

Parametric hardness? There was some nice work at SLS26

\bibliographystyle{alpha}
\bibliography{paper}

\end{document}

